Burnside 引理给出群作用下的本质不同方案数：
\[
  \operatorname{cnt}
  =\frac{1}{|G|}\sum_{g\in G}|\operatorname{Fix}(g)|.
\]
这里 $\operatorname{Fix}(g)$ 是在操作 $g$ 下保持不变的对象集合。

\subsubsection{预备知识：群}

群由一个集合 $S$ 和一个二元运算 $*$ 组成，需要满足：
\begin{enumerate}
  \item 封闭性：$\forall x,y\in S$，有 $x*y\in S$；
  \item 结合律：$x*(y*z)=(x*y)*z$；
  \item 单位元：$\exists e\in S$，使 $\forall x\in S$ 都有 $x*e=x$；
  \item 逆元：$\forall x\in S$，存在 $x^{-1}\in S$ 使 $x*x^{-1}=e$。
\end{enumerate}

\paragraph{预备例一：圆环的旋转操作}

一个圆环有 $n$ 个节点。
记 $f_i$ 为旋转 $i$ 格，操作集合为
\[
  S=\{f_0,f_1,\ldots,f_{n-1}\}.
\]
运算定义为操作复合：
\[
  f_i\circ f_j=f_{(i+j)\bmod n}.
\]
由于 $(i+j)\bmod n\in\{0,1,\ldots,n-1\}$，封闭性成立。
操作复合满足结合律，因为
\[
  (f_i\circ f_j)\circ f_k=f_i\circ(f_j\circ f_k),
\]
两边都等价于旋转 $(i+j+k)\bmod n$ 格。
单位元为 $f_0$。
任意 $f_i$ 的逆元为
\[
  f_{(n-i)\bmod n},
\]
因为二者复合为 $f_0$。
因此圆环上的所有旋转操作构成一个群。

\paragraph{预备例二：长方体的循环平移}

一个 $a\times b\times c$ 的长方体由 $abc$ 个单位小方块组成，有 $k$ 种颜色，第 $i$ 种颜色的数量为 $d_i$。
下面先证明三个方向上的所有循环平移操作构成一个群；颜色数量不影响这个操作群的证明。
每个位置记为
\[
  (x,y,z),
  \qquad
  0\le x<a,
  \quad 0\le y<b,
  \quad 0\le z<c.
\]
记 $T_{i,j,k}$ 为三个方向分别循环平移 $i,j,k$ 格，操作集合为
\[
  G=\{T_{i,j,k}\mid 0\le i<a,\ 0\le j<b,\ 0\le k<c\},
\]
所以 $|G|=abc$。
平移把位置变成
\[
  ((x+i)\bmod a,(y+j)\bmod b,(z+k)\bmod c).
\]
两个平移复合为
\[
  T_{(i_1+i_2)\bmod a,
     (j_1+j_2)\bmod b,
     (k_1+k_2)\bmod c}.
\]
复合后仍是合法平移，满足封闭性；模加法满足结合律；单位元是 $T_{0,0,0}$；逆元是
\[
  T_{(-i)\bmod a,(-j)\bmod b,(-k)\bmod c}.
\]
因此所有三维循环平移操作构成一个群。

\subsubsection{例题一：长度为 4 的黑白项链}

长度为 $4$ 的项链，每颗珠子可染成黑白两色，旋转后相同算同一种方案。
四种旋转的固定方案数依次为 $16,2,4,2$，因此
\[
  \operatorname{ans}=\frac14(16+2+4+2)=6.
\]

\subsubsection{例题二：只考虑旋转的 $k$ 色项链}

长度为 $n$，每个位置有 $k$ 种颜色。
旋转 $r$ 格会把位置分成 $\gcd(n,r)$ 个循环，固定方案数为 $k^{\gcd(n,r)}$，所以
\[
  \operatorname{ans}
  =\frac1n\sum_{r=0}^{n-1}k^{\gcd(n,r)}.
\]
按 $g=\gcd(n,r)$ 分组可得
\[
  \operatorname{ans}
  =\frac1n\sum_{g\mid n}\varphi\!\left(\frac ng\right)k^g.
\]
证明如下。
因为 $g\mid n$，两边同除以 $g$ 后有
\[
  \gcd\!\left(\frac ng,\frac rg\right)=1.
\]
对固定的 $g$，满足条件的 $r$ 的个数正是 $\varphi(n/g)$。

\subsubsection{例题三：旋转和翻转都视为相同}

长度为 $n$ 的项链有 $n$ 个旋转和 $n$ 个翻转，因此 $|G|=2n$。
旋转部分贡献
\[
  \sum_{r=0}^{n-1}k^{\gcd(n,r)}
  =\sum_{g\mid n}\varphi\!\left(\frac ng\right)k^g.
\]

当 $n$ 为奇数时，每条翻转轴经过一个珠子和一条边。
有 $1$ 个珠子固定，其余珠子两两配对，独立位置数为
\[
  1+\frac{n-1}{2}=\frac{n+1}{2}.
\]
每个翻转固定 $k^{(n+1)/2}$ 个方案，共有 $n$ 个翻转，所以
\[
  \operatorname{ans}
  =\frac1{2n}\left(
    \sum_{g\mid n}\varphi\!\left(\frac ng\right)k^g
    +nk^{(n+1)/2}
  \right).
\]

当 $n$ 为偶数时，翻转轴分为两类。
第一类经过两个相对珠子，有 $2$ 个珠子固定，其余两两配对，独立位置数为
\[
  2+\frac{n-2}{2}=\frac n2+1.
\]
这种翻转有 $n/2$ 个，贡献为
\[
  \frac n2k^{n/2+1}.
\]
第二类经过两条相对边，没有珠子固定，所有珠子两两配对，独立位置数为 $n/2$。
这种翻转有 $n/2$ 个，贡献为
\[
  \frac n2k^{n/2}.
\]
因此
\[
  \operatorname{ans}
  =\frac1{2n}\left(
    \sum_{g\mid n}\varphi\!\left(\frac ng\right)k^g
    +\frac n2k^{n/2+1}
    +\frac n2k^{n/2}
  \right).
\]

\begin{icpcwarning}[固定方案不是操作个数]
  Burnside 求和项是每个群元素固定的染色方案数。
  按 $\gcd$ 分组时，满足 $\gcd(n,r)=g$ 的旋转数是 $\varphi(n/g)$。
\end{icpcwarning}
